\chapter{Related Work -- Nghiên cứu liên quan}\label{ch:related}
Chương này phân tích những hướng liên quan theo nhiệm vụ và output contract. Danh mục phương pháp
gần đây giữ cửa sổ khảo sát đã chốt 28/10/2023--28/10/2026; tám công trình được
khảo sát có công bố đầu tiên đã xác định trong giai đoạn 2024--2026. Geo-I và các nguồn lý thuyết cũ
được dẫn như nền tảng, không được dùng để mở rộng giả tạo cửa sổ khảo sát gần đây. Danh mục có phạm
vi lựa chọn phục vụ bài toán, không phải một systematic review toàn bộ lĩnh vực.

\section{Bảo vệ truy vấn vị trí liên tục}
\subsection{TransProtect: tọa độ thay thế có xét bối cảnh}
TransProtect sử dụng mô hình chuyển động và thông tin bối cảnh để chọn tọa độ tổng hợp cho phương
tiện \cite{transprotect}. Cơ chế không chỉ nhằm tạo một điểm ở gần GPS, mà còn xem xét liệu đầu ra
có hợp lý dưới kiến thức của attacker về đường và việc di chuyển hay không. Output được xem là một
tọa độ thay thế ở mỗi mốc, phù hợp với nhiệm vụ online. Paper đánh giá sai số suy luận và tổn thất
chi phí di chuyển tới những đích có thể.

Hướng này gần với S1/S3 của luận văn vì attacker có thể kết hợp chuỗi tọa độ và context. Tuy nhiên,
travel-cost distortion của một điểm thay thế không tự tương đương Recall của danh sách POI cuối.
Khi dịch vụ trả nhiều kết quả từ nhiều query locations, cần xác định chính xác đầu ra dịch vụ và
không lấy một thành viên gần GPS nhất bằng nhãn riêng để giả lập released singleton.

Trong thực nghiệm của luận văn, hàng TransProtect* là \emph{Markov adaptation}. Nó giữ hướng chọn
miền/đích và đo chi phí phù hợp nhưng thay predictor neural của nguồn bằng proxy học từ TRAIN.
Sự thay thế này có thể ảnh hưởng chất lượng bảo vệ. Vì vậy, kết quả là đối chiếu với adaptation được
công bố đầy đủ, không được trình bày như đã tái lập đúng checkpoint và mọi chi tiết của tác giả.

\subsection{Tương quan ngữ nghĩa của chuỗi tập vị trí}
Phương pháp dummy-based xét semantic correlation của moving paths tạo các tập vị trí có thành viên
thật và các dummy \cite{semantic}. Điểm giả phải có tính nhất quán theo thời gian và ý nghĩa địa điểm,
vì một tập thoạt nhìn có nhiều thành viên vẫn có thể bị thu hẹp khi attacker kiểm tra đường đi hoặc
ngữ nghĩa. Paper sử dụng những phép đo anonymity success và chất lượng dummy trong môi trường nguồn.

Output chứa điểm thật tạo ra một bài toán đặc biệt: attacker có thể cố nhận diện thành viên đó.
Đây là khác biệt quan trọng với Q-only của GeoI-BQ, nơi GPS thật không được đưa vào tập như một thành
viên bắt buộc. Vì vậy, cùng ký hiệu K không làm hai phương pháp có cùng bảo đảm; xác suất chọn thành
viên thật và sai số định vị là hai contract khác nhau.

Mapping scope cần thận trọng. Việc xét ngữ nghĩa và liên tiếp chưa đủ để gọi là đã trực tiếp đo suy
luận điểm dừng hoặc tái dựng mọi đoạn đường trong S2/S3 của luận văn. Adaptation Semantic correlation*
dùng POI công khai và predictor thực nghiệm thay các nguồn semantic/learner gốc. Những phần chưa
tái lập được ghi riêng để người đọc không gán kết quả trên cohort này cho implementation chính thức.

\subsection{Chèn fake queries giữa các lần truy vấn thật}
Fake-query insertion bổ sung những tập hoàn toàn giả giữa các yêu cầu có vị trí thật/dummy
\cite{fake}. Mục đích là làm gián đoạn việc nối các tập theo thời gian và tăng số giải thích có thể
cho một chuỗi quan sát. Cơ chế quan tâm tới các đường khó phân biệt, sự hợp lý của fake queries và
chi phí tạo/phát chúng.

Lịch gửi là một phần của threat model. Nếu attacker biết những thời điểm mà client luôn gửi yêu cầu
thật, giả định rằng mọi packet đều có khả năng là thật không còn phù hợp. Một phép đánh giá công
bằng cần giữ tự nhiên số fake events, cho attacker nhìn tất cả sự kiện trong cùng cửa sổ và tính đủ
request/reply phát sinh. So cùng số packet trong khi bỏ packet giả sẽ thay đổi nhiệm vụ của phương
pháp, cũng như làm sai phần chi phí.

Hàng Fake-query insertion* của luận văn là road/history adaptation. Các giả thiết về đồ thị, lịch
chờ và similarity được công bố; không gọi nó là bản chính thức của paper. Kết quả S8 cũng cho thấy
đối chiếu hai tập có thể làm attacker mạnh hơn khi có một thành viên thật tương quan, nên bank kiểm
thử không nên chỉ gồm một estimator độc lập cho từng tập.

\section{Geo-I trên mạng đường và cơ chế có xét utility}
\subsection{Nền tảng Geo-I và dự đoán riêng tư}
Geo-I ràng buộc tỷ số xác suất của cùng một output giữa hai vị trí đầu vào theo khoảng cách
\cite{geoi}. Nó cung cấp một phát biểu xác suất để nối với posterior của attacker, thay vì chỉ yêu cầu
điểm giả nằm xa GPS. Các kết quả về cơ chế dự đoán riêng tư cho mobility traces cho thấy việc quản
lý quan sát mới và tái sử dụng thông tin là một phần của bài toán chuỗi \cite{predictive}.

Luận văn kế thừa hướng bảo đảm này, nhưng phải đặc tả thêm support trên mạng đường, phép kiểm tra
giữ/đổi và action chọn Q. Chứng minh của một lần thêm nhiễu không tự áp dụng cho mọi quyết định dùng
GPS trong một chương trình. Những điều kiện hậu xử lý và hợp thành được phát biểu ở chương phương
pháp, nơi có thể đối chiếu trực tiếp với giao diện input/output.

\subsection{EV querying: approximate Geo-I, dummy và xử lý cục bộ}
Cơ chế truy vấn dành cho electric vehicles kết hợp approximate Geo-I, dummy và xử lý cục bộ nhằm
giữ chất lượng lựa chọn trạm sạc \cite{ev}. Utility được đo qua cost of privacy của tác vụ tương ứng;
một tọa độ nhiễu có thể được sử dụng như bước trung gian để phục vụ quyết định của client.

Công trình này là tiền lệ quan trọng: kết hợp Geo-I, nhiều query locations và lựa chọn local không
phải bản thân một tuyên bố mới của luận văn. Điều cần làm rõ ở GeoI-BQ là policy cụ thể, cận riêng
của việc chọn đoạn, điều kiện cho toàn transcript và protocol đối chiếu giữa các dạng output.
Các giả định Edge, ngữ cảnh trạm sạc và approximate guarantee của EV querying cũng không nên được
đồng nhất với một phiên client/server có cùng clock và cap của nghiên cứu này.

\subsection{PTPPM: bảo vệ có xét mạng đường và tương quan}
Bản PTPPM được khảo sát là preprint năm 2025 về personalized trajectory protection trên road network
với spatiotemporal correlations \cite{ptppm}. Nó liên quan trực tiếp tới việc chọn miền đường,
ngân sách và sự đánh đổi distortion/utility. Những nhận định ở đây dựa trên phiên bản nguồn đã khảo
sát; luận văn không mặc định có một reproduction cùng tác vụ POI của paper.

PTPPM được đưa vào tổng hợp để định vị hướng nghiên cứu gần mô hình đề xuất. Tuy nhiên, chưa có
implementation đã audit đáp ứng contract và nguồn dữ liệu hiện tại, nên không gán điểm số của một
phương pháp khác cho hàng PTPPM. Việc không có số trong bảng thực nghiệm không làm paper mất liên
quan; nó phản ánh ranh giới fidelity của phép so sánh đang thực hiện.

\section{Danh tính, phát hành trajectory và các phạm vi liên quan}
\subsection{CPCROK và liên kết phương tiện}
CPCROK xét bảo vệ pseudonym/beacon trong mạng xe mật độ thấp bằng các cơ chế liên quan mix-zone,
thay pseudonym và fake beacons \cite{cpcrok}. Mục tiêu nối lại các quan sát của xe liên quan tới S4,
nhưng kênh VANET, danh tính pseudonym và loại bản tin khác với truy vấn POI online.

Một client gửi năm điểm Q không tự xóa account, IP, biển số hoặc định danh phần cứng. Vì vậy, luận
văn tách việc chặn thông tin bổ sung từ tọa độ khỏi một hệ thống ẩn danh danh tính hoàn chỉnh.
Dataset chỉ có family tuyến cũng không đủ để đánh giá người và phương tiện như hai target độc lập.
S4 của phần thực nghiệm bổ sung được đặt đúng tên là association proxy.

\subsection{DP-FETC và trajectory publishing offline}
DP-FETC kết hợp thống kê riêng tư, feature extraction, trajectory correlation và xử lý phân bố OD
trong bài toán phát hành dữ liệu tổng hợp \cite{dpfetc}. Các phép đo về phân bố, mutual information
hoặc độ phức tạp phản ánh mục tiêu khai thác dataset sau công bố. Chúng không trực tiếp cho biết
attacker định vị người dùng online từ một packet Q tốt đến đâu.

Quan hệ với S3/S8/S9/S10 chủ yếu là một phần về nội dung quỹ đạo, tương quan và đầu/cuối. Luận văn
không ánh xạ mọi kết quả offline thành protection đầy đủ cho các task online tương ứng. Phát hành
một tập trajectory và cung cấp dịch vụ POI liên tục có điều kiện input/output cũng như mục tiêu
utility khác nhau, nên không được ép vào cùng bảng score chỉ để đủ số phương pháp.

\subsection{SoK và yêu cầu xác định đơn vị privacy}
SoK về trajectory generation hệ thống hóa khó khăn khi kết hợp privacy và utility, nhấn mạnh việc
định nghĩa đơn vị được bảo vệ và đánh giá cả bảo đảm lẫn triển khai \cite{sok}. Nguồn này hữu ích để
phân biệt point-level, trajectory-level và những claim về người dùng hoặc nhiều phiên liên kết.

Trong luận văn, một phiên được giữ bởi token/ledger và một mô hình clock cụ thể. Bảo đảm toàn
transcript của phiên không tự là bảo đảm cho mọi phiên của một người, mọi người đồng hành hoặc
một quyết định tương lai do chính đáp án dịch vụ gây ra. Nhận định này được dùng để giới hạn scope,
không dùng SoK như một model đối chứng có số request/Recall POI tương đương.

\section{Tổng hợp, khoảng trống và hướng tiếp cận của luận văn}
Bảng~\ref{tab:related} gom các trục so sánh vào cùng một chỗ. F nghĩa là paper trực tiếp xét/đánh giá
task dưới giả định của nó; P nghĩa là một phần hoặc một setting khác; CX nghĩa là chưa đủ căn cứ
cho mapping chi tiết. Những nhãn này do luận văn ánh xạ sang S1--S10, không phải chứng chỉ bảo vệ
hoàn toàn trước mọi attacker. Scenario không được liệt kê chưa có căn cứ coverage trong audit này.
\par\begingroup\footnotesize\setstretch{1.04}\setlength{\parskip}{0pt}
\begin{longtable}{@{}P{2.75cm}P{2.15cm}P{1.6cm}P{2.4cm}P{2.5cm}P{3.194cm}@{}}
\caption{Tổng hợp nghiên cứu: cơ chế, output, scenario, metrics và fidelity}\label{tab:related}\\
\toprule
\textbf{Paper / cơ chế} & \textbf{Output} & \textbf{Scenario} & \textbf{Privacy} & \textbf{Utility / cost} & \textbf{Tái lập trong bản này} \\
\midrule\endfirsthead
\multicolumn{6}{l}{\small\textit{Bảng \thetable\ (tiếp)}}\\\toprule
\textbf{Paper / cơ chế} & \textbf{Output} & \textbf{Scenario} & \textbf{Privacy} & \textbf{Utility / cost} & \textbf{Tái lập trong bản này} \\
\midrule\endhead
\midrule\multicolumn{6}{r}{\small\textit{Còn tiếp ở trang sau}}\\\endfoot
\bottomrule\endlastfoot
\textbf{TransProtect} 2024\cite{transprotect}\newline Học chuyển động + chọn miền + obfuscation & Một tọa độ thay thế/mốc & F: S1, S3\newline P: S2 & EIE; tấn công VehiTrack & Travel-cost distortion $\Delta$c; chi phí vận hành chưa có contract chung & Chạy Markov adaptation; không có checkpoint Transformer gốc \\
\textbf{Semantic correlation} 2026\cite{semantic}\newline Tương quan ngữ nghĩa và thời gian & GPS thật + K-1 dummy & F: S1\newline P: S2, S3 & Anonymity success rate (ASR) & DER là chất lượng dummy; thời gian sinh & Chạy OSM/empirical adaptation; không tái lập AMap/LSTM \\
\textbf{Fake queries} 2026\cite{fake}\newline Chèn query toàn giả giữa query thật & Real+dummy sets và fake packets & F: S1, S3\newline P: S2 & ASR; số đường khó phân biệt & Thời gian/RSS; chưa đo Recall POI chung & Chạy road/history adaptation, giữ fake timing và mọi chi phí \\
\textbf{EV querying} 2024\cite{ev}\newline Approximate Geo-I + dummy + Edge & Vị trí nhiễu và dummy; chọn trạm sạc local & P: S1, S2, S3 & ($\varepsilon$,$\delta$)-AGeoI; IBU/Wasserstein & Cost of Privacy (CoP), tỷ lệ CoP=0 & Khảo sát; chưa có adaptation cùng dịch vụ và giả định Edge \\
\textbf{CPCROK} 2025\cite{cpcrok}\newline Mix-zone, đổi pseudonym và beacon giả & Beacon thật + tuyến beacon giả & P: S4 & Tỷ lệ khớp lại quỹ đạo $\rho$ & Số tuyến giả $\psi$; truyền thông VANET & Khác hợp đồng POI online; không đưa vào bảng score chính \\
\textbf{DP-FETC} 2025\cite{dpfetc}\newline Thống kê DP, tổng hợp trace và chỉnh OD & Bộ dữ liệu quỹ đạo offline & P: S3, S8, S9, S10 & $\varepsilon$-DP; mutual information & JSD phân bố; độ phức tạp & Offline publishing; không chấm bằng request POI online \\
\textbf{PTPPM} 2025\cite{ptppm}\newline Personalized budget, PLS, Permute-and-Flip & Vị trí đã perturb trên mạng đường & CX: S1, S3 & Geo-I + distortion privacy/EIE bound & Perturbation distance/QoS & Preprint liên quan; chưa có comparator đã audit trong repo \\
\textbf{SoK trajectory generation} 2024\cite{sok}\newline Systematization và unit of privacy & Khảo sát phát hành dữ liệu tổng hợp & Không áp dụng & Formal privacy, reconstruction/correlation threats & Mục tiêu và utility của dữ liệu tổng hợp & Nguồn bối cảnh; không là một model đối chứng online \\
\end{longtable}\endgroup\par


Ba khoảng trống dẫn tới thiết kế được nghiên cứu. Thứ nhất, cần nối tất cả thao tác dùng GPS và
quyết định chọn Q thành một cận cho transcript có trạng thái. Thứ hai, utility phải được đo qua
đáp án dịch vụ thực nhận, trong khi nhu cầu thật giữ ở local. Thứ ba, đánh giá cần thống nhất target,
cửa sổ và cost giữa các output contracts, đồng thời vẫn giữ metrics gốc khi có thể áp dụng đúng.

Một bộ metrics bổ sung vì thế không có nghĩa thay thế mọi metrics của paper. ASR của tập chứa điểm
thật đòi posterior và định nghĩa nguồn; EIE của point estimator gần với MAE định vị; DER xét chất
lượng dummy; JSD xét fidelity phân bố offline; CoP gắn với tác vụ dịch vụ cụ thể. Chương thực nghiệm
sẽ nêu cả khả năng tính và lý do N/A thay vì đổi tên các đại lượng đó thành một score chung.

GeoI-BQ được định vị như một phương pháp client-side có đặc tả và phân tích toàn transcript, với
kernel chọn đoạn có cận riêng và một luồng phục vụ nhiều nhu cầu local. Đối chiếu chính dùng ba
adaptation online phù hợp, còn các công trình chưa có reproduction cùng contract giữ vai trò
related work. Cách định vị này cho phép đánh giá một đóng góp có phạm vi rõ mà không đòi mọi paper
phải giải cùng toàn bộ mười scenario.
